\documentclass[11pt]{article}
\usepackage[margin=1in]{geometry}
\usepackage{amsmath,amssymb,amsthm}
\newcommand{\sem}[1]{[\![#1]\!]}
\usepackage{hyperref}
\usepackage[numbers,sort&compress]{natbib}

\title{Set Theoretic Estimation and Language:\\
Unification, Underspecification, and the Hyperframe Lattice}
\author{Project \texttt{thejeb}}
\date{\today}

\newtheorem{theorem}{Theorem}
\newtheorem{definition}{Definition}
\newtheorem{claim}{Claim}
\newtheorem*{conjecture}{Conjecture (outlook)}

\newcommand{\lean}[1]{\texttt{#1}}

\begin{document}
\maketitle

\begin{abstract}
This note makes an intellectual case in three movements.  First, a
historical claim: set theoretic estimation (STE) in the sense of
Combettes \citep{combettes1993} is the estimation-theoretic framing of
what two mature traditions in computational linguistics---unification
grammar and underspecification semantics---already do operationally, and
neither tradition cites Combettes; the bridge appears to be unclaimed.
Second, a structural claim, now machine-checked in the Lean module
\lean{Ste.HyperFrame}: the STE satisfaction relation is a formal context
in the sense of Formal Concept Analysis \citep{ganter1999fca}, partial
feasibility sets are exactly its lower polars, and hence the closed
feasible sets are the extents of a complete concept lattice---the
\emph{hyperframe lattice}.  This yields, as corollaries, a disproof of
naive oscillation (pure constraint accumulation is monotone) and a
rough-set \citep{pawlak1982rough} indiscernibility partition that
explains why Carlson's AEP work partition \citep{carlson2012} is
STE-invariant.  Third, an outlook: feasible sets glue over the
constraint poset in a presheaf-like way, and a sheaf-theoretic
representation is the natural next target; that direction is explicitly
conjectural.  Throughout we distinguish sharply between what is proven
in the CI-verified Lean library (cited by identifier) and what is
argued or conjectured.
\end{abstract}

\section{The thesis: an unclaimed bridge}

Combettes's foundational formulation of set theoretic estimation
\citep{combettes1993} is austere: a solution space $\Xi$, an index set
$I$ of items of information, and for each $i \in I$ a \emph{property
set} $S_i \subseteq \Xi$ of solutions consistent with that item.  The
estimate is not a point but the \emph{feasibility set}
$\bigcap_{i \in I} S_i$; acquiring information shrinks it; a point
estimate is warranted only in the \emph{ideal} case where the
intersection is a singleton.

Two traditions in computational linguistics rediscovered this schema
operationally, in different guises, without the estimation-theoretic
vocabulary and---as far as we have been able to determine---without
citing Combettes or the set-membership estimation literature at all.

\paragraph{Unification grammar.}
In unification-based approaches to grammar
\citep{shieber1986unification, carpenter1992logic}, partial information
about a linguistic object is a \emph{feature structure}, and combining
two pieces of information is \emph{unification}: the greatest lower
bound in the subsumption order, when it exists.  Each feature structure
$F$ denotes the set $\sem{F}$ of fully specified objects
it subsumes, and the defining property of unification is precisely
constraint intersection:
\[
\sem{F \sqcap G} \;=\; \sem{F} \cap \sem{G} .
\]
Unification failure is the empty intersection---in STE terms, an
infeasible (over-constrained) problem.  A parse succeeds ideally when
the accumulated constraints cut the candidate space to a singleton.
This is STE's ``intersect the property sets'' scheme executed by a
grammar, with subsumption playing the role of informativeness and
$\sqcap$ playing the role of adding an item of information.

\paragraph{Underspecification semantics.}
In underspecification semantics the denotation of an ambiguous sentence
is deliberately a \emph{set of readings}, represented compactly and
narrowed by constraints.  Minimal Recursion Semantics
\citep{copestake2005mrs} represents scope ambiguity by a flat bag of
elementary predications plus handle constraints; each way of solving
the constraints is one reading, so an MRS denotes the set of its
solutions.  Hole Semantics \citep{bos1996hole} and the Constraint
Language for Lambda Structures \citep{egg2001cls} make the logic
explicit: an underspecified representation is a constraint set, its
models are the readings, and adding a dominance constraint monotonically
prunes the model set.  Again this is exactly the STE picture: readings
are the solution space, each constraint is a property set, and the
current semantic value of the discourse is the feasibility set.

\paragraph{The claim.}
The claim is not that these fields need STE to do what they already do.
It is that STE is the common estimation-theoretic \emph{frame}: both
traditions perform set-valued estimation by constraint intersection,
with monotone refinement, over-constraint (failure) and
under-constraint (residual ambiguity) as the two failure modes, and the
ideal case as a singleton feasibility set.  Because neither literature
cites Combettes, the identification---and with it the transfer of STE's
structural results (below) to grammar and semantics---appears to be an
unclaimed bridge.  The companion note on Constraint Grammar in this
repository makes one instance of this identification fully formal; the
present note supplies the lattice-theoretic superstructure.

\section{Machine-checked: the hyperframe lattice}
\label{sec:hyperframe}

All results in this section are proven in Lean~4 over
Mathlib \citep{mathlib2020} in the module \lean{Ste.HyperFrame}, which
builds in this repository's CI; we cite each by its Lean identifier.
Fix $S : I \to \mathcal P(\Xi)$, a family of property sets.

\begin{definition}[satisfaction context; \lean{STE.sat}]
The relation $\mathrm{sat}\,S \subseteq \Xi \times I$ holds at $(a,i)$
iff $a \in S_i$.
\end{definition}

This is literally a \emph{formal context} in the sense of Formal
Concept Analysis \citep{ganter1999fca}: objects are hypotheses,
attributes are constraints, incidence is satisfaction.  The partial
feasibility set $\Phi_J = \bigcap_{i \in J} S_i$ (only the constraints
in $J \subseteq I$ enforced) then has an exact FCA description:

\begin{theorem}[the bridge; \lean{STE.partialFeasibilitySet\_eq\_lowerPolar}]
For every $J \subseteq I$, the partial feasibility set $\Phi_J$ equals
the lower polar (attribute derivation $J'$) of $J$ under
$\mathrm{sat}\,S$.
\end{theorem}

Consequently the Galois connection of the context organizes STE.  A
\emph{hyperframe} of $S$ (\lean{STE.HyperFrame}) is defined as a formal
concept of $\mathrm{sat}\,S$: a pair (extent, intent) with extent the
set of hypotheses satisfying every constraint in the intent, and intent
the set of constraints satisfied by every hypothesis in the extent.
Mathlib supplies the complete-lattice structure on concepts, so the
hyperframes of $S$ form a complete lattice with no further work.
The two directions of the correspondence are:

\begin{theorem}[feasible sets are extents;
\lean{STE.isExtent\_partialFeasibilitySet}]
Every partial feasibility set $\Phi_J$ is the extent of a hyperframe.
\end{theorem}

\begin{theorem}[extents are feasible sets;
\lean{STE.hyperFrame\_extent\_eq\_feasibilitySet}]
Every hyperframe's extent equals the feasibility set of its own
intent: $\mathrm{extent}(c) = \Phi_{\mathrm{intent}(c)}$.
\end{theorem}

So the lattice of hyperframes is exactly the lattice of
constraint-set/feasible-set fixpoints: the ``reachable'' estimation
states of the problem, closed under the Galois closure.  Enforcement is
compositional:

\begin{theorem}[superposition; \lean{STE.partialFeasibilitySet\_union}]
$\Phi_{J \cup K} = \Phi_J \cap \Phi_K$.
\end{theorem}

For the linguistic reading: the hyperframe lattice is the space of
stable ``analysis states'' of a grammar or of an underspecified
semantic representation.  Unification moves down this lattice (meet);
the maximally underspecified representation is its top; unification
failure is falling to an empty extent.  None of the linguistic claims
in Section~1 are formalized---what is formalized is the STE-side
lattice they would land in.

\section{Oscillation: the honest statement}
\label{sec:oscillation}

A recurring informal worry about iterated constraint processing is
\emph{oscillation}: as information accumulates, does the feasible set
thrash---shrinking to empty, then growing again?  The machine-checked
answer localizes the phenomenon precisely.

\begin{theorem}[no oscillation under accumulation;
\lean{STE.ncard\_partialFeasibilitySet\_antitone}]
Over a finite hypothesis space, if $J \subseteq K$ then
$|\Phi_K| \le |\Phi_J|$: accumulating constraints can only shrink the
feasible cardinality.
\end{theorem}

\begin{theorem}[oscillation requires retraction;
\lean{STE.lt\_ncard\_imp\_not\_subset}]
If $|\Phi_J| < |\Phi_K|$ then $J \not\subseteq K$: any growth of the
feasible set forces the drop (retraction) of some previously enforced
constraint.
\end{theorem}

So \emph{pure accumulation is monotone and cannot oscillate}; this
disproves the naive form of the worry.  What remains true---and this is
the honest refinement---is that oscillation is exactly a
\emph{dynamic} phenomenon of edit streams that include retraction:
reanalysis in parsing, defeasible constraints, belief revision.  The
repository's \lean{Ste.DynamicFrame} development studies that
non-monotone regime; the present result fences it off cleanly from
static STE.

\section{Rough sets and Carlson's AEP partition}
\label{sec:rough}

The context $\mathrm{sat}\,S$ also carries Pawlak's rough-set structure
\citep{pawlak1982rough}.  Define two hypotheses to be
\emph{indiscernible} (\lean{STE.Indisc}) when they lie in exactly the
same property sets.

\begin{theorem}[\lean{STE.indisc\_equivalence}]
Indiscernibility is an equivalence relation on $\Xi$.
\end{theorem}

\begin{theorem}[granularity; \lean{STE.indisc\_mem\_iff}]
Indiscernible hypotheses are co-feasible for every enforced constraint
set $J$: $a \in \Phi_J \iff b \in \Phi_J$.
\end{theorem}

The equivalence classes are the \emph{resolution floor} of the
estimation problem: no subfamily of the available constraints can
separate two indiscernible hypotheses, so every lower/upper (Pawlak)
approximation, and every feasible set, is a union of classes.  This is
the STE reading of Carlson's use of the asymptotic equipartition
property in his cryptanalytic dissertation \citep{carlson2012}: the
AEP typical-set partition used there to divide the work of key search
is an indiscernibility partition for the induced property sets, and the
granularity theorem is the formal reason the partition is STE-invariant
---constraint processing respects the classes, so work can be
partitioned along them once and for all.  (The identification of
Carlson's specific partition with \lean{STE.Indisc} for his context is
an interpretive claim of this note; what is machine-checked is the
general indiscernibility theory just cited.)

\section{Outlook: representation via sheaves}
\label{sec:sheaves}

The superposition law $\Phi_{J\cup K} = \Phi_J \cap \Phi_K$ and
antitonicity say that $J \mapsto \Phi_J$ is a limit-preserving
contravariant assignment on the poset of constraint sets: local
feasibility data over a cover $\{J_k\}$ of $J = \bigcup_k J_k$
determines feasibility over $J$, and compatible local sections should
glue.  The repository already contains a first mechanized step in this
direction, \lean{Ste.DynamicFrame.PresheafGluing}, which proves a
gluing law for feasibility data over document/constraint structure.

\begin{conjecture}
The assignment $J \mapsto \Phi_J$ extends to a sheaf on a suitable site
(the constraint poset with covers given by unions, or a Grothendieck
topology on the document poset in the dynamic setting), and the
hyperframe lattice of Section~\ref{sec:hyperframe} is recovered as its
lattice of closed subsheaf sections.  A representation theorem of this
form would subsume both the concept-lattice picture and the dynamic
gluing law.
\end{conjecture}

We flag this explicitly as \emph{not yet proven}: no sheaf condition,
site structure, or representation theorem is currently formalized, and
we deliberately do not pursue sheaf- or homotopy-theoretic
formalization here.  It is the open direction this note points at.

\bibliographystyle{plainnat}
\bibliography{../../refs}

\end{document}
